\section{Interview Guides}
\label{app:interview-guides}

\subsection{Interview Guide 1: Mini-Grid Developers/Operators}
\textbf{Introduction}

Thank you for participating. This research examines how carbon markets can support rural mini-grid development in sub-Saharan Africa. Your insights on commercial viability, implementation experiences, and policy environments are valuable for understanding barriers and opportunities.

\textbf{Background}
\begin{enumerate}
	\item Can you describe your organization's mini-grid portfolio? (Number of sites, capacity, locations, business model)
	\item What are your biggest financing challenges for mini-grid projects?
\end{enumerate}

\textbf{B. Carbon Market Engagement}
(Use the relevant subsection based on respondent profile)
For developers pursuing / having pursued carbon credits:
\begin{enumerate}[resume]
	\item What motivated your organization to pursue carbon financing for mini-grids?
	\item Walk me through your experience with the carbon certification process. What were the key steps and challenges?
	\item How significant are carbon revenues relative to your overall project economics?
	\item What transaction costs did you encounter? (Financial, time, technical capacity)
	\item Did carbon credit requirements affect how you design or operate mini-grids?
\end{enumerate}

For developers who considered but rejected carbon finance:
\begin{enumerate}[start=3]
	\item What made your organization consider carbon financing initially?
	\item What factors led to the decision not to pursue carbon credits?
	\item What would need to change to make carbon markets viable for your projects?
\end{enumerate}

For developers who never seriously considered carbon finance:
\begin{enumerate}[start=3]
	\item Are you aware of carbon market opportunities for mini-grids?
	\item Why hasn't your organization pursued carbon financing?
	\item What concerns or barriers do you perceive regarding carbon markets?
\end{enumerate}

\textbf{C. Policy and Regulatory Environment}
\begin{enumerate}[start=8]
	\item How do current regulations in [Kenya/Nigeria] support or hinder mini-grid development?
	\item Are there specific policies affecting how carbon revenues would be treated? (Taxation, tariffs, subsidies)
	\item What regulatory changes would most help scale mini-grid deployment?
\end{enumerate}

\textbf{D. Development Impacts}
\begin{enumerate}[resume]
	\item What socio-economic benefits do your mini-grids deliver to communities?
	\item How do you balance financial sustainability with maximizing community benefits?
	\item If carbon revenues were available, how might they be used? (Lower tariffs, expand capacity, community programs)
\end{enumerate}

\textbf{E. Implementation Insights}
\begin{enumerate}[resume]
	\item What factors have been critical to your successful projects?
	\item What advice would you give other developers considering carbon finance?
\end{enumerate}

\subsection{Interview Guide 2: Policymakers/Regulators}
\textbf{Introduction}

Thank you for participating. This research examines policy frameworks that could enable carbon market integration with rural mini-grids. Your perspective on regulatory environments and policy development is essential.

\textbf{A. Policy Context}
\begin{enumerate}
	\item What are your country's/region's current priorities for rural electrification?
	\item How do mini-grids fit within national energy access strategies?
	\item What financing mechanisms currently support mini-grid development?
\end{enumerate}

\textbf{B. Carbon Market Integration}
\begin{enumerate}[resume]
	\item Is carbon financing for mini-grids discussed in policy circles? If so, how?
	\item Are there existing regulations addressing how carbon revenues should be treated? (Taxation, tariff implications, licensing requirements)
	\item What coordination exists between your ministry/agency and those responsible for climate/carbon policy?
\end{enumerate}

\textbf{C. Regulatory Gaps and Barriers}
\begin{enumerate}[resume]
	\item What policy gaps currently hinder mini-grid deployment generally?
	\item What specific regulatory challenges would carbon-financed mini-grids face?
	\item How do current policies address project aggregation across multiple developers?
\end{enumerate}

\textbf{D. Enabling Environment}
\begin{enumerate}[resume]
	\item What policy changes could facilitate carbon market integration with mini-grids?
	\item What institutional capacity would be needed to implement such policies?
	\item Are there concerns about carbon financing potentially conflicting with other policy objectives? (Tariff affordability, subsidy programs, grid extension plans)
\end{enumerate}

\textbf{E. Regional Cooperation}
\begin{enumerate}[resume]
	\item What regional initiatives could support carbon-financed mini-grid development?
	\item What lessons from other countries might be applicable here?
\end{enumerate}

\textbf{F. Future Directions}
\begin{enumerate}[resume]
	\item How do you see the role of carbon markets in achieving electrification targets?
\end{enumerate}

\subsection{Interview Guide 3: Carbon Registry Official/Methodologist}
\textbf{Introduction}

Thank you for participating. This research examines barriers and opportunities for integrating carbon markets with rural mini-grids in sub-Saharan Africa. Your expertise on methodologies and certification processes is invaluable.

\textbf{A. Methodology and Standards}
\begin{enumerate}
	\item What methodologies does [Gold Standard/Verra] offer for mini-grid projects?
	\item How do these methodologies account for mini-grid characteristics? (Small scale, rural contexts, limited baseline data)
	\item What are typical certification timelines and costs for mini-grid projects?
\end{enumerate}

\textbf{B. Developer Engagement}
\begin{enumerate}[resume]
	\item What level of interest have you observed from mini-grid developers in sub-Saharan Africa?
	\item What barriers do developers most commonly face in pursuing certification?
	\item How does project scale affect certification feasibility?
\end{enumerate}

\textbf{C. Aggregation Mechanisms}
\begin{enumerate}[resume]
	\item Can you explain how Programme of Activities (PoA) or grouped project approaches work for mini-grids?
	\item What are the requirements and costs for establishing a PoA?
	\item Why has uptake of aggregation mechanisms been limited in the mini-grid sector?
\end{enumerate}

\textbf{D. Co-Benefits and Impact Measurement}
\begin{enumerate}[resume]
	\item How do current methodologies value development co-benefits beyond emission reductions? (Health, education, gender impacts)
	\item Are there discussions about enhancing co-benefit recognition in mini-grid methodologies?
\end{enumerate}

\textbf{E. Transaction Costs and Streamlining}
\begin{enumerate}[resume]
	\item What efforts are underway to reduce transaction costs for small-scale projects?
	\item Could digital monitoring or simplified verification approaches reduce costs?
\end{enumerate}

\textbf{F. Market Dynamics}
\begin{enumerate}[resume]
	\item How do you see voluntary carbon market demand for mini-grid projects?
	\item What advice would you give mini-grid developers considering carbon certification?
\end{enumerate}

\subsection{Interview Guide 4: Carbon Experts}
\textbf{Introduction}

Thank you for participating. This research examines how carbon markets can support rural mini-grid deployment in sub-Saharan Africa. Your expertise in carbon finance and project development provides valuable perspective on implementation realities.

\textbf{A. Carbon Market Landscape}
\begin{enumerate}
	\item How would you characterize the current state of carbon markets for small-scale renewable energy projects in sub-Saharan Africa?
	\item What types of renewable energy projects are successfully accessing carbon finance in the region?
\end{enumerate}

\textbf{B. Mini-Grid Specific Challenges}
\begin{enumerate}[resume]
	\item What are the primary barriers preventing mini-grid developers from accessing carbon markets?
	\item How do transaction costs for mini-grid carbon certification compare to project revenues?
	\item At what project scale do carbon economics become viable for mini-grids?
\end{enumerate}

\textbf{C. Implementation Experience}
\begin{enumerate}[resume]
	\item Have you worked with or observed mini-grid carbon projects in Africa? What were the key success factors or failure points?
	\item What role do intermediaries or aggregators play, and why aren't there more serving the mini-grid sector?
\end{enumerate}

\textbf{D. Methodological and Technical Issues}
\begin{enumerate}[resume]
	\item Are current carbon methodologies well-suited to mini-grid contexts, or do they need adaptation?
	\item What are the biggest technical challenges for mini-grids in meeting MRV requirements?
\end{enumerate}

\textbf{E. Market Access and Pricing}
\begin{enumerate}[resume]
	\item How do voluntary carbon buyers view mini-grid credits compared to other project types?
	\item What's the pricing differential?
\end{enumerate}

\subsection{Ethical Considerations}
\begin{itemize}
	\item Obtain informed consent before beginning
	\item Remind participants they can skip questions or stop anytime
	\item Ensure confidentiality and explain how data will be used
	\item Provide contact information for follow-up questions
\end{itemize}